\section{Conclusion}\label{sec:conclusion}
Private innovation intermediation and public basic access respond to different economic returns. Private fees select founders according to their willingness to pay, which in the maintained model is aligned with technical ability. A government instead compares the cost of access with direct service gains and the regional incidence of additional ventures. The conditional public-cost ceiling connects these decisions to basic and premium service values, consumer and incumbent effects, founder profit shares and cross-region matching.

The strongest analytical result is the primitive profit-share band in Proposition~1. With sufficiently weak matching and full partner participation, stronger skill complementarity can increase the value of restoring basic access even when the resulting ventures reduce regional commercialisation surplus. Positive basic-service value and privately appropriable technical returns jointly sustain this effect. The finite-matching continuous certificate confirms the stated signs within its fixed economic environment. It does not turn the conditional result into an unrestricted equilibrium theorem.

National evaluation adds the other region's welfare increment to the local public-cost ceiling. That known externality principle takes a specific form through privately selected abilities and downstream competition here. The three certified provision examples exhibit both local/national agreement and public underprovision relative to the sixteen-configuration benchmark. Their narrow cost band and the failed distributional and competition-form portability records restrict their use as policy rules. No general signed market-size or incumbent-count comparative static, optimal corrective policy, or physical-facility recommendation follows.

For evaluation, the model motivates measuring eligible firms, pre-service abilities, private fees, basic and premium use, joint-project outcomes and the regional allocation of consumer, incumbent and venture returns. These are predictions and measurement requirements, not empirical findings. The policy implication is to evaluate the direct service access lost under private pricing separately from the regional gains or losses of additional ventures, and then add foreign welfare when the decision is evaluated nationally. The primitive boundaries identify when these components point in opposing directions. This conditional applied-theory contribution does not require the configuration switch to be quantitatively large, but its limited portability prevents treating the examples as a general regional development rule.
